% =====================================================================
% RELATED WORK. Four paragraphs, from the phenomenon to the mechanism,
% ending on the work we build on directly. Describe prior work in its own
% terms; criticism only in subordinate clauses. TRAP: never say a component
% "carries" or "drives" a behaviour unless the cited work intervened.
% Citation fixes still open: Lu et al. (title/authors), Izawa et al.
% (authors), Sprejer et al. (authors). See references.bib.
% =====================================================================
\section{Related Work}
\label{sec:related}

% FLOW RW1: viability, personality is measurable in text | CHARACTER:
% personality -> language models | HAND-OFF: ends on models expressing
% personas that prompts can set
% Trimmed 2026-10-31 for the 8-page limit: four paragraphs -> three.
% FLOW RW1: personality and persona directions | HAND-OFF: where is it written
\paragraph{Personality and persona.}
The Big Five describes personality along five traits, each divided into facets \citep{soto2017next}, and personality leaves measurable traces in writing \citep{pennebaker1999linguistic, mairesse2007using}. Language models express personas that a prompt can set: traits such as sycophancy correspond to directions in a model's activations that can monitor and steer its persona \citep{chen2025persona}, and steering vectors more generally shift what a model generates \citep{subramani2022extracting}. These methods act on the residual stream as a whole, which leaves open where in the network a persona is written.

% FLOW RW2: mirroring and drift | HAND-OFF: fixed positions, fixed content
\paragraph{Mirroring and drift.}
Over long conversations, assistants depart from their default persona, and clamping activations along a single assistant direction reduces the departure \citep{lu2026assistant}. Assistants also shift towards the user's personality: their activations move along trait components after users interact with them, and \citet{sprejer2026white} leave open whether this reflects the model's own persona or its representation of the user's style. Inside the model the boundary between user and assistant roles is weak, and style can dominate role markers \citep{ye2026prompt}, so probes trained across positions can confuse whose style they detect. We therefore read activations at fixed positions and hold content fixed.

% FLOW RW3: persona in components | HAND-OFF: the question our dataset answers
\paragraph{Persona in specific components.}
Activation patching has located persona-driven reasoning in specific components of a model \citep{poonia2025dissecting}. Most directly, \citet{izawa2026steering} find that a sparse set of attention heads, which they call style modulation heads, writes an assistant's style, and validate them with system prompts and steering. Their personas are assigned through the system prompt and are not Big Five traits, and they do not test a style that arrives through the user's message.
